\chapter{Tests et validation}

\section*{Introduction}
\addcontentsline{toc}{section}{Introduction}

La plateforme décrite au chapitre précédent doit maintenant faire la preuve de son bon
fonctionnement. Ce chapitre commence par une série de tests ciblés sur les mécanismes de
sécurité et de résilience : signature HMAC, déduplication, circuit breaker et gardien
disque. Il présente ensuite les résultats obtenus sur des \textit{Pull Requests}
réelles, dont la revue de bout en bout de la PR de démonstration, avant de conclure sur
la validation globale de la plateforme, de la checklist de mise en production au
durcissement de la configuration et à l'analyse de menace du moteur de revue.

% ─────────────────────────────────────────────────────────────────────────────
\section{Tests de validation}

    Pour vérifier le bon fonctionnement de chaque mécanisme de sécurité et de résilience
    de la solution, nous avons conduit une série de tests unitaires sur les composants
    critiques de la plateforme. Ces tests reproduisent des conditions réelles d'erreur ou
    d'attaque pour valider les comportements attendus.

    \subsection{Test de la validation HMAC-SHA256}

        Le premier test vérifie que l'agent rejette correctement les webhooks dont la
        signature est invalide ou absente. La figure \ref{fig:test_hmac} présente le
        résultat du test : un webhook avec une signature incorrecte reçoit un code HTTP
        403 Forbidden, tandis qu'un webhook correctement signé est accepté avec un code 202.

        \begin{figure}[H]
            \centering
            \frame{\includegraphics[width=0.88\columnwidth]{test_hmac_validation}}
            \caption{Test de la validation HMAC-SHA256 — rejet 403 d'une signature invalide}
            \label{fig:test_hmac}
        \end{figure}

    \subsection{Test de la déduplication Redis}

        Le deuxième test vérifie que deux webhooks identiques (même dépôt, même numéro
        de PR, même SHA de commit) ne déclenchent qu'un seul pipeline. La figure
        \ref{fig:test_dedup} montre que le second webhook retourne un code 202 mais que
        l'agent l'ignore silencieusement grâce à la clé de déduplication Redis
        \texttt{SET NX} ; le premier pipeline se termine alors sans interruption.

        \begin{figure}[H]
            \centering
            \frame{\includegraphics[width=0.88\columnwidth]{test_deduplication}}
            \caption{Test de la déduplication Redis — un seul pipeline pour deux webhooks identiques}
            \label{fig:test_dedup}
        \end{figure}

    \subsection{Test du circuit breaker LLM}

        Le troisième test vérifie le comportement de l'agent lorsque Ollama est
        indisponible. Après trois échecs consécutifs d'appels LLM, le circuit breaker
        passe en état OPEN et active le fallback de classification regex. La figure
        \ref{fig:test_circuit} montre que le pipeline continue de fonctionner avec une
        classification basée sur les extensions de fichiers, sans bloquer sur
        l'indisponibilité du modèle.

        \begin{figure}[H]
            \centering
            \frame{\includegraphics[width=0.88\columnwidth]{test_circuit_breaker}}
            \caption{Test du circuit breaker LLM — bascule sur le fallback de classification}
            \label{fig:test_circuit}
        \end{figure}

    \subsection{Test du gardien disque}

        Le quatrième test vérifie le déclenchement automatique du nettoyage Docker
        lorsque l'utilisation du disque dépasse 90\%. En simulant un seuil critique via
        la jauge Prometheus \texttt{agent\_\allowbreak{}disk\_\allowbreak{}used\_\allowbreak{}percent}, nous avons vérifié que
        la commande \texttt{docker builder prune -f} est exécutée automatiquement et
        que le résultat est transmis sur Slack, comme le montre la figure
        \ref{fig:test_disk_guard}.

        \begin{figure}[H]
            \centering
            \frame{\includegraphics[width=0.75\columnwidth]{test_disk_guard}}
            \caption{Test du gardien disque — nettoyage automatique au-delà du seuil critique}
            \label{fig:test_disk_guard}
        \end{figure}

% ─────────────────────────────────────────────────────────────────────────────
\section{Résultats}

    Cette section rassemble les résultats obtenus une fois la plateforme en service. Elle
    part d'une revue de démonstration détaillée, l'élargit à un corpus de
    \textit{Pull Requests} réelles, la confronte à la revue manuelle, puis présente les
    vérifications de conformité, d'observabilité et de durcissement de la plateforme.

    \subsection{Résultat de la revue sur une Pull Request de démonstration}

        Le pipeline a été validé de bout en bout sur la PR~\#18 du dépôt
        \texttt{GhaiethFerchichi/\allowbreak{}Vunl-application}, ouverte le 28~mai~2026 pour
        un module bancaire de démonstration contenant cinq vulnérabilités types.
        Les figures~\ref{fig:pr_full_result}, \ref{fig:pr_full_result_2}
        et~\ref{fig:pr_full_result_3} en donnent la vue d'ensemble, scindée en
        trois parties pour la lisibilité : commentaires de revue, suggestions
        \textit{inline}, verdict \texttt{BLOCK} et échec du check
        \texttt{security-ai-agent}.

        \begin{figure}[H]
            \centering
            \frame{\includegraphics[trim=0 5109 0 0, clip, width=0.82\columnwidth]{pr_review_full_result}}
            \caption{PR~\#18 — vue d'ensemble de la revue sur GitHub (1/3) :
                     en-tête de la \textit{Pull Request} et premiers commentaires de revue}
            \label{fig:pr_full_result}
        \end{figure}

        \begin{figure}[H]
            \centering
            \frame{\includegraphics[trim=0 2554 0 2555, clip, width=0.82\columnwidth]{pr_review_full_result}}
            \caption{PR~\#18 — vue d'ensemble de la revue sur GitHub (2/3) :
                     suite des commentaires \textit{inline} sur le diff}
            \label{fig:pr_full_result_2}
        \end{figure}

        \begin{figure}[H]
            \centering
            \frame{\includegraphics[trim=0 0 0 5109, clip, width=0.82\columnwidth]{pr_review_full_result}}
            \caption{PR~\#18 — vue d'ensemble de la revue sur GitHub (3/3) :
                     revue de sécurité consolidée et verdict \texttt{BLOCK}}
            \label{fig:pr_full_result_3}
        \end{figure}

        Le tableau~\ref{tab:pr_results} synthétise le résultat de cette revue.

        \begin{table}[H]
        \centering
        \begin{tabular}{|c|c|c|c|c|c|}
        \hline
        \textbf{PR} & \textbf{Risque} & \textbf{Verdict} & \textbf{Classification} & \textbf{Durée} & \textbf{Commentaires} \\
        \hline
        \#18 & HIGH & BLOCK & feature & \textasciitilde15 min & 5 inline \\
        \hline
        \end{tabular}
        \caption[Résultat de la revue PR de démonstration]{Résultat de la revue PR de démonstration — dépôt GhaiethFerchichi/Vunl-application}
        \label{tab:pr_results}
        \end{table}

        Le fichier de démonstration contenait cinq vulnérabilités volontairement
        introduites, ce qui fournit une vérité terrain explicite. Le tableau
        \ref{tab:pr18_detection} confronte ces cinq failles aux \textit{findings}
        effectivement produits par l'agent.

        \begin{table}[H]
        \centering
        \begin{tabular}{|m{5.2cm}|c|c|m{3.2cm}|}
        \hline
        \textbf{Vulnérabilité plantée} & \textbf{Ligne} & \textbf{Détectée} & \textbf{Source} \\
        \hline
        Injection SQL & 6 & \checkmark & Analyse LLM \\
        \hline
        Secrets en dur (clés AWS, mot de passe BD) & 11--13 & \checkmark & Analyse LLM \\
        \hline
        Path traversal & 17 & \checkmark & LLM + Semgrep \\
        \hline
        Injection de commande & 22 & \checkmark & LLM + Semgrep \\
        \hline
        Hachage faible (MD5) & 27 & \checkmark & Analyse LLM \\
        \hline
        \end{tabular}
        \caption{Détection sur la PR~\#18 : cinq vulnérabilités plantées, cinq détectées}
        \label{tab:pr18_detection}
        \end{table}

        Sur ce jeu contrôlé, l'agent a détecté les cinq vulnérabilités plantées (rappel
        de 5/5), sans \textit{finding} erroné dans le fichier modifié. Deux observations
        en découlent. D'abord, les secrets codés
        en dur (lignes 11 à 13) ont été repérés par le modèle alors que Gitleaks ne les
        avait pas signalés : le LLM et les scanners spécialisés se complètent plutôt
        qu'ils ne se recouvrent. Ensuite, ces chiffres portent sur des vulnérabilités
        introduites à dessein et de facture classique ; ils valident la chaîne de
        détection de bout en bout, sans préjuger du comportement de l'agent face à des
        failles plus subtiles ou inédites.

    \subsection{Revues produites sur un corpus de Pull Requests}

        La base de connaissances a accumulé huit revues réelles au fil des sprints,
        toutes conservées dans la table \texttt{pr\_\allowbreak{}reviews}. La figure
        \ref{fig:pr_reviews_corpus} en présente les cinq revues de PRs applicatives
        (classification \textit{feature}), celles pour lesquelles le moteur Semgrep entre
        en jeu : pour chacune, le score de risque, le verdict rendu et le nombre de
        \textit{findings} remontés par Semgrep.

        \begin{figure}[H]
            \centering
            \frame{\includegraphics[width=\textwidth]{pr_reviews_corpus}}
            \caption{Revues de PRs applicatives — risque, verdict et \textit{findings} SAST (table \texttt{pr\_reviews})}
            \label{fig:pr_reviews_corpus}
        \end{figure}

        Semgrep a relevé entre 28 et 37 \textit{findings} de niveau ERROR par PR, et le
        verdict suit le risque : \texttt{REQUEST\_\allowbreak{}CHANGES} pour les deux PRs classées
        \textsc{medium}, \texttt{BLOCK} pour les trois classées \textsc{high}. La PR~\#18,
        détaillée plus haut, en est le cas le plus net : cinq vulnérabilités plantées,
        toutes détectées, verdict \texttt{BLOCK}. Sur CPU, l'inférence du modèle 14B (environ
        3 tokens par seconde) reste le principal facteur de durée d'une revue ; c'est là que
        la migration GPU évoquée en conclusion changerait l'échelle.

    \subsection{Comparaison avec la revue manuelle}

        Le tableau \ref{tab:baseline} met en regard le processus manuel décrit au
        chapitre~1 et la chaîne automatisée mise en place.

        \begin{table}[H]
        \centering
        \begin{tabular}{|m{3.2cm}|m{4.7cm}|m{4.7cm}|}
        \hline
        \textbf{Critère} & \textbf{Revue manuelle} & \textbf{BTE Security AI Agent} \\
        \hline
        Délai de revue & Jusqu'à 24~heures, selon la disponibilité et l'expertise du relecteur & 15 à 25 minutes selon la PR \\
        \hline
        Couverture & Variable selon la disponibilité et l'expertise du relecteur & Cinq scanners et analyse LLM, OWASP Top 10 systématique \\
        \hline
        Gate CI/CD & Aucun blocage technique & Verdict \texttt{APPROVE} / \texttt{REQUEST\_\allowbreak{}CHANGES} / \texttt{BLOCK} \\
        \hline
        Secrets et dépendances & Souvent manqués à l'œil nu & Gitleaks et OSV-Scanner systématiques \\
        \hline
        Disponibilité & Tributaire du relecteur du jour & Automatique, à chaque \textit{Pull Request} \\
        \hline
        Traçabilité & Fil de discussion GitHub & Base PostgreSQL et tableaux de bord \\
        \hline
        \end{tabular}
        \caption{Revue de sécurité manuelle vs automatisée}
        \label{tab:baseline}
        \end{table}

        L'automatisation ne remplace pas le relecteur humain : elle lui retire le travail
        répétitif de détection systématique et le recentre sur les décisions
        d'architecture et sur les cas que l'agent remonte en \texttt{BLOCK}.

    \subsection{Résultats de l'analyse de conformité et observabilité}

        La figure \ref{fig:prom_alerts_ok} présente l'état des quinze règles d'alerte
        Prometheus, toutes inactives, ce qui confirme l'absence d'incident en production lors
        de la validation finale.

        \begin{figure}[H]
            \centering
            \frame{\includegraphics[width=0.88\columnwidth]{prometheus_alerts_inactive}}
            \caption{Prometheus — 15 règles d'alerte inactives}
            \label{fig:prom_alerts_ok}
        \end{figure}

    \subsection{Validation de la plateforme}

        Le tableau \ref{tab:validation} présente la checklist de validation complète
        effectuée sur la plateforme avant la présentation finale.

        \begin{table}[H]
        \centering
        \begin{tabular}{|m{6cm}|m{5cm}|c|}
        \hline
        \textbf{Vérification} & \textbf{Méthode} & \textbf{Statut} \\
        \hline
        Pipeline PR end-to-end & PR réelle → commentaire en 15–25 min & \checkmark \\
        \hline
        5 scanners exécutés & Artefacts dans ./artifacts/scans/ & \checkmark \\
        \hline
        LLM local (aucune fuite) & Port 11434 interne uniquement & \checkmark \\
        \hline
        4 cibles Prometheus vertes & /prometheus/targets & \checkmark \\
        \hline
        15 règles d'alerte inactives & /prometheus/alerts & \checkmark \\
        \hline
        AlertManager → Slack fonctionnel & Test alert → message Slack reçu & \checkmark \\
        \hline
        3 dashboards Grafana chargés & Données réelles dans chaque panneau & \checkmark \\
        \hline
        VictoriaMetrics données 90j & Requête Grafana > 30 jours & \checkmark \\
        \hline
        Anti-hallucination chat validé & "CPU usage?" → appel outil + valeur réelle & \checkmark \\
        \hline
        Gardien disque actif & Jauge agent\_disk\_used\_percent Prometheus & \checkmark \\
        \hline
        Health digest Slack quotidien & Message reçu à 09h00 UTC & \checkmark \\
        \hline
        Métriques host node-exporter & PromQL node\_memory\_MemAvailable & \checkmark \\
        \hline
        nginx Basic Auth /ui protégé & Accès refusé sans credentials & \checkmark \\
        \hline
        Base de données PR reviews & SELECT pr\_reviews — 8 revues (PR~\#11 à \#18) & \checkmark \\
        \hline
        \end{tabular}
        \caption{Checklist de validation complète de la plateforme}
        \label{tab:validation}
        \end{table}

    \subsection{Durcissement de la plateforme et analyse de menace}

        La checklist précédente couvre le périmètre fonctionnel et l'observabilité. En
        prolongement de l'audit du 28 avril, une revue de durcissement a été conduite
        sur la configuration de la plateforme afin d'aligner l'exploitation sur les
        principes de défense en profondeur attendus d'un environnement bancaire. Le
        tableau \ref{tab:hardening_gaps} en synthétise les axes.

        \begin{table}[H]
        \centering
        \begin{tabular}{|m{3.4cm}|m{4cm}|m{5.2cm}|}
        \hline
        \textbf{Axe} & \textbf{Principe} & \textbf{Mise en œuvre} \\
        \hline
        Cloisonnement réseau & Un point d'entrée unique & nginx seul exposé en frontal (80/443) ; les services internes communiquent sur le réseau Docker interne, le diagnostic passe par \texttt{docker exec} ou par une route nginx authentifiée \\
        \hline
        Authentification & Défense en profondeur & Double couche sur les interfaces : HTTP Basic Auth au niveau de nginx, puis comptes applicatifs natifs (Grafana, interface de chat) ; mot de passe dédié sur le cache Redis ; création des comptes du chat réservée à l'administrateur \\
        \hline
        Confidentialité de l'inférence & Aucune donnée ne quitte le VPS & Modèles servis par Ollama sur le réseau interne uniquement, aucun appel vers une API d'inférence externe \\
        \hline
        Intégrité des événements & Authentification des sources & Signature HMAC-SHA256 vérifiée sur chaque webhook GitHub, rejet en 403 de toute requête non signée \\
        \hline
        \end{tabular}
        \caption{Axes de durcissement de la configuration de production}
        \label{tab:hardening_gaps}
        \end{table}

        Trois chantiers de consolidation prolongent cette démarche. L'automatisation
        du renouvellement des certificats TLS via Let's Encrypt simplifiera la
        gestion du chiffrement des échanges entre les utilisateurs et la plateforme.
        Une sauvegarde quotidienne de PostgreSQL, qui héberge à la fois la base de
        connaissances sécurité et les \textit{checkpoints} LangGraph, protégera contre
        la perte du volume Docker ; la panne VictoriaMetrics d'avril a rappelé qu'un
        tel incident peut survenir sans prévenir. Enfin, la généralisation des
        \textit{healthchecks} Docker à l'ensemble des conteneurs permettra au démon
        Docker de détecter et de relancer un service défaillant de lui-même, sans
        attendre que les règles Prometheus finissent par déclencher une alerte.

        Le durcissement ne se limite pas à la surface réseau : une analyse de menace a
        également été menée sur le moteur de revue lui-même. Le diff soumis au modèle
        14B est rédigé par l'auteur de la \textit{Pull Request} analysée ; il constitue
        à ce titre une entrée non fiable, et un développeur malveillant pourrait y
        glisser des instructions destinées à détourner la revue, par exemple un
        commentaire de code enjoignant au modèle d'émettre un verdict \texttt{APPROVE}.
        L'architecture oppose deux barrières à ce scénario. D'une part, les cinq
        scanners sont déterministes et s'exécutent indépendamment du modèle : un secret
        ou une CVE relevé par un scanner est enregistré dans PostgreSQL quel que soit
        le texte produit par le LLM, et ne disparaît pas si celui-ci est manipulé.
        D'autre part, la température fixée à 0,1 et l'extraction structurée du bloc
        JSON réduisent la latitude laissée au modèle dans la formulation de son
        verdict. Le plan de durcissement prévoit d'aller plus loin sur ce point :
        isoler explicitement le diff par des délimiteurs dans le \textit{prompt}, et
        faire reposer le \textit{gate} de blocage sur la sévérité des \textit{findings}
        déterministes des scanners plutôt que sur le seul verdict du modèle, de sorte
        qu'aucune injection dans le code revu ne puisse relâcher le contrôle de
        sécurité.

\section*{Conclusion}
\addcontentsline{toc}{section}{Conclusion}

Ce chapitre a soumis la plateforme à l'épreuve des faits. Les tests des
mécanismes critiques ont confirmé les comportements attendus : rejet en 403 des webhooks
non signés, déduplication des notifications répétées, bascule du circuit breaker sur le
fallback de classification et déclenchement automatique du gardien disque. Sur la PR de
démonstration~\#18, l'agent a détecté les cinq vulnérabilités plantées et rendu un
verdict \texttt{BLOCK} en une quinzaine de minutes, commentaires \textit{inline} et
suggestions de correction publiés directement sur GitHub ; le corpus de huit revues
réelles confirme la cohérence des verdicts avec les scores de risque. Rapportée au
processus manuel décrit au chapitre~1, la chaîne automatisée ramène le délai de revue de
près de vingt-quatre heures à moins d'une demi-heure, avec une couverture systématique.
La checklist de validation, les axes de durcissement et l'analyse de menace du moteur de
revue complètent ce bilan ; la conclusion générale en dresse maintenant la synthèse.
